\section{Introducción: de Long CoT a Interleaved Thinking}

Desde que OpenAI presentó o1 (Reasoning Model) en septiembre de 2024, y con el lanzamiento de DeepSeek R1 en enero de 2025, el «Thinking Model» se convirtió en un paradigma importante en el campo de los grandes modelos. Sin embargo, el Long Chain-of-Thought (Long CoT) tradicional presenta deficiencias evidentes.

K2 Thinking propone un paradigma completamente nuevo------\textbf{Interleaved Thinking} (razonamiento intercalado/pensamiento intercalado), orientado a escenarios Agentic, que mediante aprendizaje por refuerzo entrena al modelo para alternar entre pensar y responder.

\begin{importantbox}{Problemas del Long CoT tradicional}
\begin{itemize}
  \item Antes de la answer final hay un proceso extenso de long-CoT
  \item La latencia hasta el primer token (Time to First Token, TTFT) es sumamente alta, y la experiencia de usuario es deficiente
  \item Es difícil definir un reward de grano fino, y la supervisión del proceso resulta complicada
\end{itemize}
\end{importantbox}

